\documentclass{article}
\usepackage{amsmath}
\usepackage{listings}
\usepackage{geometry}

\geometry{a4paper, margin=0.7in}

\title{\vspace{-1.5cm}Homework 5}
\author{Showkot Hossain}
\date{}

\begin{document}

\maketitle
\vspace{-0.5cm}

\section*{Part 1: Conceptual Questions}
\vspace{-0.2cm}

\subsection*{Q1: }
\vspace{-0.1cm}
% \texttt{<html>.*?<head>.*?<title>.*?</title>.*?</head>.*?<body[\^{}>]*>.*?</body>.*?</html>}
No, this regular expression is not safe. The problems are:

\begin{enumerate}
    \setlength\itemsep{-0.1em}
    \item \textbf{ReDoS Vulnerability:} The combination of non-greedy quantifiers (\texttt{.*?}) and the character class in \texttt{<body[\^{}>]*>} can cause catastrophic backtracking. If input contains many \texttt{<body} sequences without closing \texttt{>}, the regex engine may take exponential time, causing Denial of Service.
    
    \item \textbf{Dot/Newline Issue:} The dot (\texttt{.}) doesn't match newlines by default in many engines. Attackers can bypass validation by inserting newlines, or valid multiline HTML may fail to match.
    
    \item \textbf{HTML Parsing Limitations:} Regular expressions cannot properly parse HTML (a context-free language). This regex is easily fooled by nested tags, comments (\texttt{<!-- </html> -->}), or tag-like strings, making it unsafe for security decisions.
\end{enumerate}

\subsection*{Q2: }
To reach \texttt{printf("everywhere\textbackslash n")}, the following conditions must be satisfied:

\begin{enumerate}
    \setlength\itemsep{-0.1em}
    \item First \texttt{if (x > 5)} is true: $x > 5$
    \item Second \texttt{if (y > 7)} is false (enter else): $y \le 7$ (i.e., $\neg(y > 7)$)
    \item Third \texttt{if (x < 20)} is true: $x < 20$
\end{enumerate}

\vspace{0.1cm}
\noindent The path condition is: \boxed{(x > 5) \land (y \le 7) \land (x < 20)}

\end{document}
